\documentclass{article}
\usepackage{amsmath, amssymb, amsthm, geometry, hyperref, mathrsfs, mathtools, tikz}
\usepackage[colorlinks=true, linkcolor=blue]{hyperref}
\geometry{margin=1in}

\title{Comprehensive Practice Exam II: Advanced Proof Techniques and Logic}
\author{Advanced Mathematics Simulator}
\date{\today}

\begin{document}

\maketitle
\tableofcontents
\newpage

\section{Concept Checklist \& Problem Characterization}
This practice exam rigorously covers the following problem archetypes based on the provided logic and proof theory framework. Each question set contains a "Similar" (standard), "Harder" (advanced algebraic manipulation), and "Twist" (conceptual leap) variant.
\begin{itemize}
    \item \textbf{Type 1: Mathematical Induction (PMI):} Direct algebraic manipulation of sums and divisibility.
    \item \textbf{Type 2: Modular Arithmetic \& Finite Rings:} Modular tables, inverses, parity arguments as a filter.
    \item \textbf{Type 3: Impossibility Proofs (Irrationality):} Reductio ad absurdum, prime factorization logic, rationality bounds.
    \item \textbf{Type 4: Contrapositive \& Contradiction on Parity:} Negation of implications, algebraic parity substitution.
    \item \textbf{Type 5: Non-Constructive Existence (IVT):} Continuous function bounds, uniqueness proofs via algebraic inequality contradictions.
    \item \textbf{Type 6: Divisibility Logic \& Direct Arguments:} Biconditional splitting, logical translations, algebraic isolation of integers.
\end{itemize}

\vspace{1cm}

\section{Problem Statements}

\subsection{Question Set 1: Mathematical Induction}
\noindent\textbf{Problem 1 (Similar):} Use the Principle of Mathematical Induction to prove that for all $n \ge 1$:
\[ \sum_{i=1}^n (2i - 1) = n^2 \]

\noindent\textbf{Problem 2 (Harder):} Use the Principle of Mathematical Induction to prove that for all $n \ge 1$:
\[ \sum_{i=1}^n \frac{1}{i(i+1)} = \frac{n}{n+1} \]

\noindent\textbf{Problem 3 (Twist):} Use the Principle of Mathematical Induction to prove the divisibility property for all $n \ge 0$:
\[ 5 \mid (7^n - 2^n) \]

\subsection{Question Set 2: Modular Arithmetic}
\noindent\textbf{Problem 4 (Similar):} 
(a) Construct the multiplication table for $\mathbb{Z}_6$. Identify all zero-divisors. \\
(b) Let $x, y \in \mathbb{Z}$. Prove via the Direct Method: If $x \equiv 2 \pmod 3$ and $y \equiv 1 \pmod 3$, then $3 \mid (x^2 + 2y)$.

\noindent\textbf{Problem 5 (Harder):} 
(a) Construct the multiplication table for $\mathbb{Z}_7$. What is the multiplicative inverse of $[4]$? \\
(b) Let $a, b \in \mathbb{Z}$. Prove via the Direct Method: If $a \equiv 3 \pmod 4$ and $b \equiv 1 \pmod 4$, then $8 \mid (a^2 - b^2)$.

\noindent\textbf{Problem 6 (Twist):} 
Prove by exhaustive cases (using the Division Algorithm) that for any integer $n$, the square of $n$ satisfies either:
\[ n^2 \equiv 0 \pmod 4 \quad \text{or} \quad n^2 \equiv 1 \pmod 4 \]

\subsection{Question Set 3: Irrationality Proofs}
\noindent\textbf{Problem 7 (Similar):} 
(a) Prove by contradiction that $\sqrt{5} \in \mathbb{I}$ (is irrational). \\
(b) Give an example of an irrational number strictly between $7$ and $8$.

\noindent\textbf{Problem 8 (Harder):} 
(a) Let $p$ be any prime number. Prove by contradiction that $\sqrt{p} \in \mathbb{I}$. \\
(b) Give an example of an irrational number strictly between $\sqrt{5}$ and $\sqrt{7}$.

\noindent\textbf{Problem 9 (Twist):} 
Prove by contradiction that the product of any non-zero rational number $q \in \mathbb{Q} \setminus \{0\}$ and any irrational number $x \in \mathbb{I}$ is always an irrational number.

\subsection{Question Set 4: Contrapositive and Contradiction Logic}
\noindent\textbf{Problem 10 (Similar):} Define statement $S: \forall a,b \in \mathbb{Z}, a \text{ even and } b \text{ odd} \implies 4 \nmid (a^2 + 2b^2)$. \\
(i) Write the logical negation, $\sim S$. \\
(ii) Prove $S$ is true using a proof by contradiction.

\noindent\textbf{Problem 11 (Harder):} Define statement $S: \forall x,y \in \mathbb{Z}, 3 \nmid x \text{ and } 3 \nmid y \implies 3 \mid (x^2 + y^2 + 1)$. \\
(i) Write the logical negation, $\sim S$. \\
(ii) Prove $S$ is true using a proof by contradiction.

\noindent\textbf{Problem 12 (Twist):} Define statement $S: \forall n \in \mathbb{Z}, \text{ if } n^3 + 5 \text{ is odd, then } n \text{ is even}$. \\
(i) Write the contrapositive of $S$. \\
(ii) Prove $S$ is true using a proof by contrapositive.

\subsection{Question Set 5: Intermediate Value Theorem and Uniqueness}
\noindent\textbf{Problem 13 (Similar):} 
(a) State the Intermediate Value Theorem. \\
(b) Prove that the polynomial $f(x) = x^3 + x - 1 = 0$ has at least one real root in the interval $[0,1]$. \\
(c) Assume $x_0$ and $x_1$ are both roots in $[0,1]$. If $x_0 < x_1$, arrive at an algebraic contradiction to prove the root is unique.

\noindent\textbf{Problem 14 (Harder):} 
(a) Prove that the polynomial $f(x) = x^5 + 3x^3 + x - 2 = 0$ has at least one real root in $[0,1]$. \\
(b) Assume $x_0 < x_1$ and both are roots in the interval. Derive an algebraic contradiction to show uniqueness.

\noindent\textbf{Problem 15 (Twist):} 
Prove that the polynomial $f(x) = 2x^7 + x^5 + 3 = 0$ has exactly one real root. (Hint: Prove existence on an arbitrary interval like $[-2, 0]$, then prove uniqueness globally).

\subsection{Question Set 6: Divisibility Logic}
\noindent\textbf{Problem 16 (Similar):} Define $S: \forall x,y \in \mathbb{Z}, 7 \mid x \text{ and } 7 \mid (x+y) \implies 7 \mid y$. \\
(i) Write the contrapositive. (ii) Write the negation. \\
(iii) Prove $S$ directly. (iv) Prove $S$ by contradiction.

\noindent\textbf{Problem 17 (Harder):} Define $S: \forall a,b,c \in \mathbb{Z}, a \mid b \text{ and } a \mid c \implies a \mid (3b - 2c)$. \\
(i) Write the contrapositive. (ii) Write the negation. \\
(iii) Prove $S$ directly. (iv) Prove $S$ by contradiction.

\noindent\textbf{Problem 18 (Twist):} Define $S: \forall a,b,n,m \in \mathbb{Z}, \text{ if } a \equiv b \pmod n \text{ and } m \mid n, \text{ then } a \equiv b \pmod m$. \\
(i) Write the contrapositive. (ii) Write the negation. \\
(iii) Prove $S$ directly. (iv) Prove $S$ by contradiction.

\newpage
\section{Detailed Solutions}

\subsection{Solutions to Set 1}
\begin{proof}[Solution to Problem 1]
Let $P(n)$ be the statement $\sum_{i=1}^n (2i - 1) = n^2$. \\
\textbf{Base Case ($n=1$):} $\sum_{i=1}^1 (2i - 1) = 2(1) - 1 = 1$. The RHS evaluates to $1^2 = 1$. Thus $P(1)$ holds. \\
\textbf{Inductive Step:} Assume $P(k)$ is true for some $k \ge 1$. We must prove $P(k+1)$.
\[ \sum_{i=1}^{k+1} (2i - 1) = \left( \sum_{i=1}^k (2i - 1) \right) + (2(k+1) - 1) \]
By the inductive hypothesis:
\[ = k^2 + (2k + 2 - 1) = k^2 + 2k + 1 \]
Factoring the perfect square:
\[ = (k+1)^2 \]
This is exactly $P(k+1)$. By PMI, $P(n)$ is true for all $n \ge 1$.
\end{proof}

\begin{proof}[Solution to Problem 2]
Let $P(n): \sum_{i=1}^n \frac{1}{i(i+1)} = \frac{n}{n+1}$. \\
\textbf{Base Case ($n=1$):} LHS $= \frac{1}{1(2)} = \frac{1}{2}$. RHS $= \frac{1}{1+1} = \frac{1}{2}$. $P(1)$ holds. \\
\textbf{Inductive Step:} Assume $P(k)$ is true for $k \ge 1$. We evaluate the series at $k+1$:
\[ \sum_{i=1}^{k+1} \frac{1}{i(i+1)} = \left( \sum_{i=1}^k \frac{1}{i(i+1)} \right) + \frac{1}{(k+1)(k+2)} \]
By the inductive hypothesis:
\[ = \frac{k}{k+1} + \frac{1}{(k+1)(k+2)} \]
Find a common denominator:
\[ = \frac{k(k+2) + 1}{(k+1)(k+2)} = \frac{k^2 + 2k + 1}{(k+1)(k+2)} = \frac{(k+1)^2}{(k+1)(k+2)} = \frac{k+1}{k+2} \]
This matches $P(k+1)$. By PMI, the statement holds for all $n \ge 1$.
\end{proof}

\begin{proof}[Solution to Problem 3]
Let $P(n)$ be the statement that $5 \mid (7^n - 2^n)$. \\
\textbf{Base Case ($n=0$):} $7^0 - 2^0 = 1 - 1 = 0$. Since $5 \cdot 0 = 0$, $5 \mid 0$. $P(0)$ holds. \\
\textbf{Inductive Step:} Assume $P(k)$ is true for some $k \ge 0$. Then $7^k - 2^k = 5m$ for some integer $m$. We evaluate $7^{k+1} - 2^{k+1}$:
\[ 7^{k+1} - 2^{k+1} = 7 \cdot 7^k - 2 \cdot 2^k \]
We can rewrite $7 \cdot 7^k$ as $(5 + 2)7^k = 5 \cdot 7^k + 2 \cdot 7^k$:
\[ = 5 \cdot 7^k + 2 \cdot 7^k - 2 \cdot 2^k = 5 \cdot 7^k + 2(7^k - 2^k) \]
Substitute $7^k - 2^k = 5m$ from the inductive hypothesis:
\[ = 5 \cdot 7^k + 2(5m) = 5(7^k + 2m) \]
Since $(7^k + 2m) \in \mathbb{Z}$, $5 \mid (7^{k+1} - 2^{k+1})$. By PMI, $P(n)$ holds $\forall n \ge 0$.
\end{proof}

\subsection{Solutions to Set 2}
\begin{proof}[Solution to Problem 4]
\textbf{(a)} Table for $\mathbb{Z}_6$:
\begin{center}
$
\begin{array}{c|cccccc}
\cdot & [0] & [1] & [2] & [3] & [4] & [5] \\
\hline
[0] & [0] & [0] & [0] & [0] & [0] & [0] \\
[1] & [0] & [1] & [2] & [3] & [4] & [5] \\
[2] & [0] & [2] & [4] & [0] & [2] & [4] \\
[3] & [0] & [3] & [0] & [3] & [0] & [3] \\
[4] & [0] & [4] & [2] & [0] & [4] & [2] \\
[5] & [0] & [5] & [4] & [3] & [2] & [1]
\end{array}
$
\end{center}
The zero divisors are $[2], [3], \text{ and } [4]$ because they multiply by non-zero elements to yield $[0]$. \\
\textbf{(b)} Assume $x \equiv 2 \pmod 3$ and $y \equiv 1 \pmod 3$. By definition, $x = 3q + 2$ and $y = 3p + 1$ for $p, q \in \mathbb{Z}$.
\[ x^2 + 2y = (3q+2)^2 + 2(3p+1) = 9q^2 + 12q + 4 + 6p + 2 = 9q^2 + 12q + 6p + 6 \]
\[ = 3(3q^2 + 4q + 2p + 2) \]
Since the inner term is an integer, $3 \mid (x^2+2y)$.
\end{proof}

\begin{proof}[Solution to Problem 5]
\textbf{(a)} Table for $\mathbb{Z}_7$:
\begin{center}
$
\begin{array}{c|ccccccc}
\cdot & [0] & [1] & [2] & [3] & [4] & [5] & [6] \\
\hline
[0] & [0] & [0] & [0] & [0] & [0] & [0] & [0] \\
[1] & [0] & [1] & [2] & [3] & [4] & [5] & [6] \\
[2] & [0] & [2] & [4] & [6] & [1] & [3] & [5] \\
[3] & [0] & [3] & [6] & [2] & [5] & [1] & [4] \\
[4] & [0] & [4] & [1] & [5] & [2] & [6] & [3] \\
[5] & [0] & [5] & [3] & [1] & [6] & [4] & [2] \\
[6] & [0] & [6] & [5] & [4] & [3] & [2] & [1]
\end{array}
$
\end{center}
The multiplicative inverse of $[4]$ is $[2]$, since $[4] \cdot [2] = [8] = [1]$. \\
\textbf{(b)} Assume $a \equiv 3 \pmod 4$ and $b \equiv 1 \pmod 4$. Then $a = 4q+3$ and $b = 4p+1$ for $q, p \in \mathbb{Z}$.
\[ a^2 - b^2 = (4q+3)^2 - (4p+1)^2 = (16q^2 + 24q + 9) - (16p^2 + 8p + 1) \]
\[ = 16q^2 - 16p^2 + 24q - 8p + 8 = 8(2q^2 - 2p^2 + 3q - p + 1) \]
Since the inner term is an integer, $8 \mid (a^2-b^2)$.
\end{proof}

\begin{proof}[Solution to Problem 6]
By the Division Algorithm, any integer $n$ can be expressed as $n = 2k$ (even) or $n = 2k+1$ (odd) for some $k \in \mathbb{Z}$. We evaluate $n^2$ in both cases:
\begin{itemize}
    \item \textbf{Case 1 ($n = 2k$):} $n^2 = (2k)^2 = 4k^2$. Since $4k^2 = 4(k^2) + 0$, $n^2 \equiv 0 \pmod 4$.
    \item \textbf{Case 2 ($n = 2k+1$):} $n^2 = (2k+1)^2 = 4k^2 + 4k + 1 = 4(k^2+k) + 1$. Since $k^2+k \in \mathbb{Z}$, $n^2 \equiv 1 \pmod 4$.
\end{itemize}
In all mutually exhaustive cases, $n^2 \equiv 0 \pmod 4$ or $n^2 \equiv 1 \pmod 4$. 
\end{proof}

\subsection{Solutions to Set 3}
\begin{proof}[Solution to Problem 7]
\textbf{(a)} Assume for contradiction that $\sqrt{5} \in \mathbb{Q}$. Then $\sqrt{5} = p/q$ where $p,q \in \mathbb{Z}$, $q \neq 0$, and $\gcd(p,q)=1$. Squaring both sides yields $5 = p^2/q^2 \implies 5q^2 = p^2$. This means $5 \mid p^2$. Since $5$ is prime, Euclid's Lemma dictates $5 \mid p$. Thus $p = 5k$ for $k \in \mathbb{Z}$. Substitute this back: $5q^2 = (5k)^2 = 25k^2 \implies q^2 = 5k^2$. This implies $5 \mid q^2$, and thus $5 \mid q$. Both $p$ and $q$ share a factor of $5$, contradicting $\gcd(p,q)=1$. Thus $\sqrt{5} \in \mathbb{I}$. \\
\textbf{(b)} An irrational number between $7$ and $8$ is $7 + \frac{\sqrt{2}}{2}$.
\end{proof}

\begin{proof}[Solution to Problem 8]
\textbf{(a)} Assume $\sqrt{p} = a/b$ in lowest terms ($\gcd(a,b)=1$). $pb^2 = a^2$. This implies $p \mid a^2$. Since $p$ is prime, Euclid's Lemma states $p \mid a$. Let $a = pk$ for $k \in \mathbb{Z}$. Substitute back: $pb^2 = (pk)^2 = p^2k^2$. Dividing by $p$ yields $b^2 = pk^2$. This implies $p \mid b^2$, and since $p$ is prime, $p \mid b$. Both $a$ and $b$ are divisible by $p$, contradicting that they are in lowest terms. Thus $\sqrt{p} \in \mathbb{I}$. \\
\textbf{(b)} The arithmetic mean of two distinct real numbers lies strictly between them. Thus, $\frac{\sqrt{5} + \sqrt{7}}{2}$ is an irrational number strictly between $\sqrt{5}$ and $\sqrt{7}$.
\end{proof}

\begin{proof}[Solution to Problem 9]
Assume for contradiction that the product is rational. Let $q = a/b \in \mathbb{Q}$ ($a, b \neq 0$) and $x \in \mathbb{I}$. Assume $qx = r$, where $r = c/d \in \mathbb{Q}$ ($d \neq 0$). We isolate $x$:
\[ x = \frac{r}{q} = \frac{c/d}{a/b} = \frac{bc}{ad} \]
Since integers are closed under multiplication, $bc \in \mathbb{Z}$ and $ad \in \mathbb{Z}$. Since $a, d \neq 0$, $ad \neq 0$. This expresses $x$ as a ratio of integers, meaning $x \in \mathbb{Q}$. This contradicts the premise that $x \in \mathbb{I}$. Therefore, $qx$ must be irrational.
\end{proof}

\subsection{Solutions to Set 4}
\begin{proof}[Solution to Problem 10]
\textbf{(i)} $\sim S: \exists a, b \in \mathbb{Z}, a \text{ even and } b \text{ odd and } 4 \mid (a^2+2b^2)$. \\
\textbf{(ii)} Assume $\sim S$ is true. Then $a=2k$ and $b=2m+1$ for integers $k,m$, and $a^2+2b^2 = 4j$ for some integer $j$. We evaluate $a^2+2b^2$:
\[ a^2+2b^2 = (2k)^2 + 2(2m+1)^2 = 4k^2 + 2(4m^2+4m+1) = 4k^2 + 8m^2 + 8m + 2 = 4(k^2 + 2m^2 + 2m) + 2 \]
This implies $a^2+2b^2$ leaves a remainder of $2$ when divided by $4$. However, our assumption states it equals $4j$ (remainder $0$). This contradiction implies $S$ is true.
\end{proof}

\begin{proof}[Solution to Problem 11]
\textbf{(i)} $\sim S: \exists x,y \in \mathbb{Z}, 3 \nmid x \text{ and } 3 \nmid y \text{ and } 3 \nmid (x^2+y^2+1)$. \\
\textbf{(ii)} Assume $\sim S$ is true. Since $3 \nmid x$, $x \equiv 1 \text{ or } 2 \pmod 3$. In either case, $x^2 \equiv 1 \pmod 3$. Similarly, since $3 \nmid y$, $y^2 \equiv 1 \pmod 3$. 
Evaluating the sum:
\[ x^2 + y^2 + 1 \equiv 1 + 1 + 1 \equiv 3 \equiv 0 \pmod 3 \]
This mathematically implies $3 \mid (x^2+y^2+1)$. However, our negated premise assumes $3 \nmid (x^2+y^2+1)$. This is a direct contradiction. Thus $S$ is true.
\end{proof}

\begin{proof}[Solution to Problem 12]
\textbf{(i)} Contrapositive: $\forall n \in \mathbb{Z}, \text{ if } n \text{ is odd, then } n^3 + 5 \text{ is even}$. \\
\textbf{(ii)} Assume $n$ is odd. Then $n = 2k+1$ for $k \in \mathbb{Z}$. Substitute into the expression:
\[ (2k+1)^3 + 5 = (8k^3 + 12k^2 + 6k + 1) + 5 = 8k^3 + 12k^2 + 6k + 6 = 2(4k^3 + 6k^2 + 3k + 3) \]
Since $(4k^3 + 6k^2 + 3k + 3) \in \mathbb{Z}$, the expression is of the form $2m$, which is even. Because the contrapositive is true, the original statement $S$ is true.
\end{proof}

\subsection{Solutions to Set 5}
\begin{proof}[Solution to Problem 13]
\textbf{(a)} \textit{Intermediate Value Theorem:} Let $f$ be continuous on the closed interval $[a,b]$. If $k$ is any number between $f(a)$ and $f(b)$, there exists at least one number $c \in [a,b]$ such that $f(c) = k$. \\
\textbf{(b)} Let $f(x) = x^3+x-1$. $f$ is a polynomial, thus continuous on $[0,1]$. $f(0) = -1 < 0$. $f(1) = 1+1-1 = 1 > 0$. By IVT, since $0$ is between $-1$ and $1$, $\exists x_0 \in [0,1]$ such that $f(x_0)=0$. \\
\textbf{(c)} Assume $x_0 < x_1$ and $f(x_0) = 0 = f(x_1)$. Then $(x_0^3+x_0-1) - (x_1^3+x_1-1) = 0$. Grouping terms: $(x_0^3 - x_1^3) + (x_0 - x_1) = 0$. Since $x_0 < x_1$, we know $x_0 - x_1 < 0$. Additionally, $x_0^3 < x_1^3$, so $x_0^3 - x_1^3 < 0$. The sum of two strictly negative numbers must be $<0$, but the equation states it equals $0$. Contradiction.
\end{proof}

\begin{proof}[Solution to Problem 14]
\textbf{(a)} Let $f(x) = x^5+3x^3+x-2$. Continuous on $[0,1]$. $f(0) = -2 < 0$. $f(1) = 1+3+1-2 = 3 > 0$. By IVT, $\exists x_0 \in [0,1]$ such that $f(x_0)=0$. \\
\textbf{(b)} Assume $x_0 < x_1$ and $f(x_0) = 0 = f(x_1)$. Subtracting the two yields:
\[ (x_0^5 - x_1^5) + 3(x_0^3 - x_1^3) + (x_0 - x_1) = 0 \]
Since $x_0 < x_1$, $x_0^k < x_1^k$ for all odd $k \ge 1$. Therefore, $(x_0^5 - x_1^5) < 0$, $3(x_0^3 - x_1^3) < 0$, and $(x_0 - x_1) < 0$. A sum of strictly negative numbers cannot equal zero. Contradiction.
\end{proof}

\begin{proof}[Solution to Problem 15]
Let $f(x) = 2x^7+x^5+3$. $f$ is continuous on $\mathbb{R}$. $f(-2) = 2(-128) - 32 + 3 = -285 < 0$. $f(0) = 3 > 0$. By IVT, there exists at least one root $x_0 \in [-2,0]$. \\
Assume $f(x)$ has two real roots, $x_0$ and $x_1$, with $x_0 < x_1$. Then $f(x_0) = f(x_1) = 0$.
\[ (2x_0^7 + x_0^5 + 3) - (2x_1^7 + x_1^5 + 3) = 0 \implies 2(x_0^7 - x_1^7) + (x_0^5 - x_1^5) = 0 \]
Since $x_0 < x_1$, it is a known property of odd powers that $x_0^k < x_1^k$, meaning $(x_0^7 - x_1^7) < 0$ and $(x_0^5 - x_1^5) < 0$. 
The equation requires the sum of two strictly negative real numbers to be $0$. This is a contradiction. Thus $f(x)$ has exactly one real root.
\end{proof}

\subsection{Solutions to Set 6}
\begin{proof}[Solution to Problem 16]
\textbf{(i)} Contrapositive: $\forall x,y \in \mathbb{Z}, 7 \nmid y \implies 7 \nmid x \text{ or } 7 \nmid (x+y)$. \\
\textbf{(ii)} Negation: $\exists x,y \in \mathbb{Z}, 7 \mid x \text{ and } 7 \mid (x+y) \text{ and } 7 \nmid y$. \\
\textbf{(iii) Direct Proof:} Assume $7 \mid x$ and $7 \mid (x+y)$. Then $x = 7k$ and $x+y = 7m$ for integers $k,m$. Substitute $x$: $7k + y = 7m \implies y = 7m - 7k = 7(m-k)$. Since $m-k \in \mathbb{Z}$, $7 \mid y$. \\
\textbf{(iv) Contradiction Proof:} Assume the negation is true. $x = 7k$, $x+y = 7m$, and $7 \nmid y$. Then $y = (x+y) - x = 7m - 7k = 7(m-k)$. This implies $7 \mid y$, which contradicts our assumption that $7 \nmid y$.
\end{proof}

\begin{proof}[Solution to Problem 17]
\textbf{(i)} Contrapositive: $\forall a,b,c \in \mathbb{Z}, a \nmid (3b-2c) \implies a \nmid b \text{ or } a \nmid c$. \\
\textbf{(ii)} Negation: $\exists a,b,c \in \mathbb{Z}, a \mid b \text{ and } a \mid c \text{ and } a \nmid (3b-2c)$. \\
\textbf{(iii) Direct Proof:} Assume $a \mid b$ and $a \mid c$. $b = ak$ and $c = am$ for integers $k,m$. Substitute into the target expression: $3b - 2c = 3(ak) - 2(am) = a(3k - 2m)$. Since $3k-2m \in \mathbb{Z}$, $a \mid (3b-2c)$. \\
\textbf{(iv) Contradiction Proof:} Assume $a \mid b$, $a \mid c$, but $a \nmid (3b-2c)$. Since $a \mid b$ and $a \mid c$, $3b-2c = a(3k-2m)$ as shown above. This definitively proves $a \mid (3b-2c)$, directly contradicting $a \nmid (3b-2c)$.
\end{proof}

\begin{proof}[Solution to Problem 18]
\textbf{(i)} Contrapositive: $\forall a,b,n,m \in \mathbb{Z}, a \not\equiv b \pmod m \implies a \not\equiv b \pmod n \text{ or } m \nmid n$. \\
\textbf{(ii)} Negation: $\exists a,b,n,m \in \mathbb{Z}, a \equiv b \pmod n \text{ and } m \mid n \text{ and } a \not\equiv b \pmod m$. \\
\textbf{(iii) Direct Proof:} Assume $a \equiv b \pmod n$ and $m \mid n$. By definition, $a-b = nk$ and $n = mp$ for integers $k,p$. Substitute $n$: $a-b = (mp)k = m(pk)$. Since $pk \in \mathbb{Z}$, $m \mid (a-b)$, which means $a \equiv b \pmod m$. \\
\textbf{(iv) Contradiction Proof:} Assume $a \equiv b \pmod n$, $m \mid n$, and $a \not\equiv b \pmod m$. From the first two assumptions, $a-b = m(pk)$ as shown above. This implies $m \mid (a-b)$, meaning $a \equiv b \pmod m$. This contradicts the assumption $a \not\equiv b \pmod m$.
\end{proof}

\newpage
\section{Concept Checklist \& Question Mapping}
\begin{itemize}
    \item \textbf{Type 1: Mathematical Induction (PMI):} Problems 1, 2, 3.
    \item \textbf{Type 2: Modular Arithmetic \& Finite Rings:} Problems 4, 5, 6.
    \item \textbf{Type 3: Impossibility Proofs (Irrationality):} Problems 7, 8, 9.
    \item \textbf{Type 4: Contrapositive \& Contradiction on Parity:} Problems 10, 11, 12.
    \item \textbf{Type 5: Non-Constructive Existence (IVT):} Problems 13, 14, 15.
    \item \textbf{Type 6: Divisibility Logic \& Direct Arguments:} Problems 16, 17, 18.
\end{itemize}

\vspace{1cm}
\begin{center}
\begin{tikzpicture}[
    node distance=2cm and 3cm,
    box/.style={draw, rectangle, rounded corners, align=center, minimum width=3cm, minimum height=1cm},
    arrow/.style={->, thick}
]

% Nodes
\node[box, fill=green!10] (goal) {Goal: Prove $P \implies Q$};
\node[box, fill=blue!10, below left=of goal] (direct) {Assume $P$ \\ $\vdots$ \\ Deduce $Q$};
\node[box, fill=orange!10, below=of goal] (contra) {Assume $\sim Q$ \\ $\vdots$ \\ Deduce $\sim P$};
\node[box, fill=red!10, below right=of goal] (contradict) {Assume $P \land \sim Q$ \\ $\vdots$ \\ Deduce Contradiction ($\bot$)};

% Arrows
\draw[arrow] (goal) -- (direct) node[midway, above left] {Direct};
\draw[arrow] (goal) -- (contra) node[midway, right] {Contrapositive};
\draw[arrow] (goal) -- (contradict) node[midway, above right] {Contradiction};

\end{tikzpicture}
\end{center}

\end{document}